\chapter{文件系统实现}

\section{文件系统结构}
\concept{用户视图与系统视图}
\begin{points}
  \item 用户视图中，文件是带名字的逻辑数据集合，通过系统调用访问。
  \item 系统视图中，文件是磁盘数据块的集合，文件系统负责从逻辑文件偏移映射到物理磁盘块。
  \item 文件系统基本操作单位是数据块 block，而磁盘物理基本单位是扇区 sector。
  \item 操作系统通常一次读写一个块，块可由多个连续扇区组成。
\end{points}
\plain{把这个知识点放回本章主线理解：它回答的是 OS 为了管理资源、隐藏硬件、保证并发正确性而采用的某个对象、规则或步骤。}
\exam{常考选择/填空/简答，让你判断概念归属、比较相近概念，或按“定义-作用-机制-易错点”展开。}


\concept{分层结构}
\begin{points}
  \item 逻辑文件系统：处理文件名、目录、权限、FCB/inode 等。
  \item 文件组织模块：完成逻辑块到物理块的映射。
  \item 基本文件系统：向设备驱动发出读写块请求。
  \item I/O 控制层：设备驱动和中断处理程序控制硬件。
\end{points}
\plain{把这个知识点放回本章主线理解：它回答的是 OS 为了管理资源、隐藏硬件、保证并发正确性而采用的某个对象、规则或步骤。}
\exam{常考选择/填空/简答，让你判断概念归属、比较相近概念，或按“定义-作用-机制-易错点”展开。}


\concept{FCB 与 inode}
\begin{points}
  \item FCB 文件控制块保存文件属性和存储位置信息。
  \item inode 是类 Unix 系统中的索引节点，保存文件类型、权限、大小、时间、数据块地址等元数据。
  \item 目录项通常保存文件名和 inode 号，文件名不一定存放在 inode 内。
  \item inode 与文件物理分配方式有关，因为 inode 中需要保存指向数据块或索引块的指针。
\end{points}
\plain{FCB/inode 保存文件属性和物理位置，是文件系统找到文件内容的核心索引。}
\exam{常考目录项是否要放全部 FCB、inode 与文件名分离的好处。}


\section{目录实现}
\concept{线性列表目录}
\begin{points}
  \item 目录项按线性表保存，查找时顺序扫描。
  \item 优点：实现简单。
  \item 缺点：目录项多时查找慢，平均查找约需扫描一半目录项。
  \item 目录项访问次数题常用“目录项数/每块目录项数/平均扫描块数”计算。
\end{points}
\plain{目录就是文件名到文件控制信息的映射表，路径就是沿目录树找到文件。}
\exam{常考一级/两级/树形目录优缺点、平均查找磁盘访问次数。}


\concept{哈希目录}
\begin{points}
  \item 根据文件名计算哈希值，快速定位目录项。
  \item 优点：查找速度快。
  \item 缺点：需要处理哈希冲突，目录扩展和哈希表维护更复杂。
\end{points}
\plain{目录就是文件名到文件控制信息的映射表，路径就是沿目录树找到文件。}
\exam{常考一级/两级/树形目录优缺点、平均查找磁盘访问次数。}


\section{文件存储空间分配方法}
\concept{连续分配}
\begin{points}
  \item 文件占用磁盘上一组连续块，目录项只需记录起始块号和长度。
  \item 优点：顺序访问速度快，随机访问也容易，寻道少。
  \item 缺点：需要预先知道文件大小，文件增长困难，会产生外部碎片。
  \item 适合大小较稳定、顺序访问多的文件。
\end{points}
\plain{连续分配要求一整块内存，容易产生碎片；分页把连续要求打碎。}
\exam{常考内部/外部碎片区别、首次/最佳/最坏适应算法、伙伴系统合并。}


\concept{链式分配}
\begin{points}
  \item 文件由多个磁盘块组成链表，每个块包含指向下一个块的指针。
  \item 优点：无外部碎片，文件容易增长。
  \item 缺点：随机访问差，访问第 $k$ 个块必须从第一个块顺链找；指针损坏会影响后续数据。
  \item FAT 是链式分配的改进，把链接指针集中存放在文件分配表中。
\end{points}
\plain{每个块指向下一个块，不怕文件增长，但找第 n 块必须顺着链走。}
\exam{常考链式分配不适合随机访问，FAT 如何把指针集中到表中改善。}


\concept{索引分配}
\begin{points}
  \item 为每个文件建立索引块，索引块中保存该文件所有数据块号。
  \item 目录项保存索引块地址。
  \item 优点：支持直接访问，无外部碎片，文件可离散存放。
  \item 缺点：索引块本身占空间；小文件可能浪费索引块，大文件需要多级索引。
\end{points}
\plain{把文件所有数据块号集中放在索引块里，既支持离散分配，又方便随机访问。}
\exam{常考单级/多级索引最大文件长度计算：每个索引块能放多少块号，再乘数据块大小。}


\concept{多级索引与最大文件长度}
\begin{points}
  \item 若块大小为 $B$ 字节，每个块号占 $p$ 字节，则一个索引块可保存 $B/p$ 个块号。
  \item 一级索引最大可索引 $(B/p)$ 个数据块，最大文件大小为 $(B/p)\times B$。
  \item 二级索引最大可索引 $(B/p)^2$ 个数据块，最大文件大小为 $(B/p)^2\times B$。
  \item 三级索引同理为 $(B/p)^3\times B$。
\end{points}
\plain{把文件所有数据块号集中放在索引块里，既支持离散分配，又方便随机访问。}
\exam{常考单级/多级索引最大文件长度计算：每个索引块能放多少块号，再乘数据块大小。}

\section{空闲空间管理}
\concept{位示图/位图}
\begin{points}
  \item 用一个二进制位表示一个磁盘块是否空闲。
  \item 优点：查找连续空闲块方便，结构紧凑。
  \item 缺点：磁盘很大时位图本身也较大，需要常驻或缓存。
\end{points}
\plain{文件系统要知道哪些盘块空着，常用位图、链表、成组链接等方法管理。}
\exam{常考位示图位号与盘块号换算、空闲链表/成组链接特点。}


\concept{空闲链表}
\begin{points}
  \item 把所有空闲块链接成链表。
  \item 优点：释放和分配单个块较简单。
  \item 缺点：查找连续空闲块效率低，链指针损坏可能影响管理。
\end{points}
\plain{文件系统要知道哪些盘块空着，常用位图、链表、成组链接等方法管理。}
\exam{常考位示图位号与盘块号换算、空闲链表/成组链接特点。}


\concept{成组链接法}
\begin{points}
  \item 把空闲块分组，每组第一块记录下一组空闲块号。
  \item 兼顾链表和批量管理，是 Unix 早期常用方法。
  \item 适合大量空闲块的快速分配和回收。
\end{points}
\plain{硬链接共享同一个 inode，符号链接是另一个文件，内容是目标路径。打开文件表保存运行时读写状态。}
\exam{常考硬/软链接区别、引用计数、系统级/进程级打开文件表。}


\concept{计数法}
\begin{points}
  \item 用“起始空闲块号 + 连续空闲块数”表示一段连续空闲区。
  \item 适合磁盘空闲块成片连续的情况。
  \item 比逐块链表更紧凑，但碎片多时记录数量增加。
\end{points}
\plain{把这个知识点放回本章主线理解：它回答的是 OS 为了管理资源、隐藏硬件、保证并发正确性而采用的某个对象、规则或步骤。}
\exam{常考选择/填空/简答，让你判断概念归属、比较相近概念，或按“定义-作用-机制-易错点”展开。}


\section{文件共享实现}
\concept{基于索引节点的共享}
\begin{points}
  \item 多个目录项指向同一个 inode，实现硬链接共享。
  \item inode 中维护引用计数，记录有多少目录项引用它。
  \item 删除目录项只减少引用计数，不一定删除文件内容。
\end{points}
\plain{多个进程都想用同一批资源，所以必须规定能不能一起用、谁先用、用完怎么释放。}
\exam{常考互斥共享/同时共享举例，或者问并发和共享为什么是 OS 最基本特征。}


\concept{基于符号链接的共享}
\begin{points}
  \item 符号链接文件保存目标路径名。
  \item 访问符号链接时，系统根据路径再去查找目标文件。
  \item 优点是灵活、可跨文件系统；缺点是访问多一次路径解析，目标删除后链接失效。
\end{points}
\plain{多个进程都想用同一批资源，所以必须规定能不能一起用、谁先用、用完怎么释放。}
\exam{常考互斥共享/同时共享举例，或者问并发和共享为什么是 OS 最基本特征。}


\chapter{大容量存储系统}
